\section{Preliminaries}

\subsection{Segala's Framework}

\begin{definition}[Probabilistic Labelled Transition System]
    \lean{PLTS}
    %\leanok

    A probabilistic labelled transition system (PLTS) is a tuple $\pcal := (Q,q_\star,L,T)$
    where $Q$ is a (countable) set of states,
    $q_\star \in Q$ is the initial state,
    $L$ is a (countable) set of labels and
    $T \subseteq Q \times (L \cup \{\tau\}) \times \dcal(Q)$ is a set of probabilistic transitions.
\end{definition}

A partial execution of $\pcal$ is a sequence
$q_0.l_0.q_1.l_1\cdots \in (Q.L)^*.Q \cup (Q.L)^\omega$ such that, for every $n \in \nat$, there
exists $\mu_n \in \dcal(Q)$ with $(q_n,l_n,\mu_n) \in T$ and $q_{n+1} \in \bar{\mu_n}$. The set of
finite executions is denoted $\ecal^*$. A partial execution is called an execution if
$q_0 = q_\star$. The trace of a partial execution is its projection onto $L$. The set of traces of
$\pcal$ is denoted $\tcal$. A scheduler $s$ is a function from $\ecal^*$ to $\dcal(T \cup \{\bot\})$
such that, for all $e \in \ecal^*$ and $t \in \bar{s(e)}$, $t = \bot$ or $src(t) = last(e)$. The
set of schedulers is denoted $Sch$.

\begin{definition}[Partial Probabilistic Execution]
    \lean{Partial Probabilistic Execution}
    %\leanok

    Given $\mu \in \dcal(Q)$ and $s \in Sch$, one can define the partial probabilistic execution
    induced by $(\mu,s)$ as the function $\efrak_{\mu,s}$ from $\ecal^*$ to $[0,1]$ such that if
    $e \in Q$ then $\efrak_{\mu,s}(e) = \mu(e)$; otherwise, $e = e'.l.q$ and
    $\efrak_{\mu,s}(e) =
    \efrak_{\mu,s}(e')\cdot\sum_{t \in T \mid lbl(t) = l}s(e')(t)\cdot out(t)(q)$.
\end{definition}

A partial probabilistic execution is called a probabilistic execution if it is induced by a pair
$(\mu,s)$ such that $\mu = \delta_{q_\star}$. If $s \in Sch$ terminates almost surely from $\mu$
(i.e. $\sum_{e \in \ecal^*} \efrak_{\mu,s}(e)\cdot s(e)(\bot) = 1$) then
$\qfrak_{\mu,s} = last^{-1}\circ\efrak_{\mu,s}$ is a well-defined probability distribution over
states. We denote $\mu \xRightarrow{l} \mu'$ if and only if there exists $s \in Sch$ such that
$\delta_l = \tfrak_{\mu,s}$ and $\mu' = \qfrak_{\mu,s}$. Partial probabilistic executions induce a
measure on the $\sigma$-field induced by cones of partial executions, and the trace function is
measurable from the $\sigma$-field induced by cones of partial executions to the one induced by
cones of traces. Thus, $\tfrak_{\mu,s} = tr^{-1}\circ\efrak_{\mu,s}$ is well-defined and called a
trace distribution. The set of trace distributions achievable by $\pcal$ is defined as those
induced by a probabilistic execution of $\pcal$ and denoted $ \text{tdist}(\pcal)$.
A hyper-transition of a PLTS $\pcal$ is a triple $(\mu,l,\mu')$ such that, for all $q \in \bar\mu$,
there exists $(q,l,\mu_q) \in T$ such that $\mu' = \sum_{q \in Q}\mu(q)\cdot\mu_q$.
The set of hyper-transitions is denoted $H$.
The trace distribution pre-order $\pcal \subseteq^t \pcal'$ is defined by
$ \text{tdist}(\pcal) \subseteq  \text{tdist}(\pcal')$.

\begin{definition}[Parallel Composition of PLTS]
    \lean{Parallel Composition}
    %\leanok

    Given two PLTSs $\pcal$ and $\pcal'$ and $S \subseteq L \cap L'$, the composed PLTS
    $\pcal \parallel_S \pcal'$ is defined as
    $(Q \times Q', (q_\star, q_\star'), L \cup L', T \parallel_S T')$ where
    $((q,q'), l, \mu'') \in T \parallel_S T'$ if and only if
    \begin{itemize}
        \item $l \in S$ and there exist $(q,l,\mu) \in T$ and $(q',l,\mu') \in T'$ such that
        $\mu \otimes \mu' = \mu''$; or
        \item $l \not\in S$ and there exists $(q,l,\mu) \in T$ such that
        $\mu \otimes \delta_{q'} = \mu''$; or
        \item $l \not\in S$ and there exists $(q',l,\mu') \in T'$ such that
        $\delta_q \otimes \mu' = \mu''$.
    \end{itemize}
\end{definition}

The trace distribution pre-order is not stable under $\parallel_S$: $\pcal \subseteq^t \pcal'$ does
not imply $\pcal \parallel_S \pcal'' \subseteq^t \pcal' \parallel_S \pcal''$
\cite{Segala1995Compositional}.
Therefore, the largest pre-order included in $\subseteq^t$ stable under $\parallel_S$ is considered:
$\pcal \sqsubseteq^t \pcal'$ if and only if, for all PLTSs $\pcal''$ and
$S \subseteq L \cap L' \cap L''$,
$\pcal \parallel_S \pcal'' \subseteq^t \pcal' \parallel_S \pcal''$.

\begin{lemma}
    \label{lem:precong}
    If $\pcal \sqsubseteq^t \pcal'$ then, for all PLTSs $\pcal''$ and
    $S \subseteq L \cap L' \cap L''$,
    $\pcal \parallel_S \pcal'' \sqsubseteq^t \pcal' \parallel_S \pcal''$
\end{lemma}
\begin{proof}
    This is a direct consequence of the associativity of $\parallel_S$.
\end{proof}

\begin{definition}[Abstraction Operator]
    Given a PLTS $\pcal$ and $I \subseteq L$, the abstract PLTS $\tau_I(\pcal)$ is obtained
    by replacing all labels in $I$ by $\tau$.
\end{definition}

\begin{lemma}
    \label{lem:abstract}
    If $\pcal \sqsubseteq^t \pcal'$ then, for all $I \subseteq L$,
    $\tau_I(\pcal) \sqsubseteq^t \tau_I(\pcal')$.
\end{lemma}

Given $R \subseteq A \times \dcal(B)$, the probabilistic lifting of $R$, denoted $\dcal(R)$, is
defined by $(\mu_A,\mu_B) \in \dcal(R)$ if and only if there exists a witness $w : R \rightarrow 
[0,1]$ such that $\mu_A = a \mapsto \sum_{\mu \in \dcal(B)}w(a,\mu)$ and $\mu_B = \sum_{(a, \mu) 
\in R}w(a,\mu)\cdot\mu$.

\begin{definition}[Probabilistic Forward Simulation]
    \lean{Probabilistic Forward Simulation}
    %\leanok

    Given two PLTSs $\pcal$ and $\pcal'$, a probabilistic forward simulation of $\pcal$ by $\pcal'$
    is a relation $\scal \subseteq Q \times \dcal(Q')$ such that
    $(q_\star,\delta_{q_\star'}) \in \scal$ and, for all $(q,\mu') \in \scal$ and $(q,l,\mu) \in T$,
    there exists $\mu' \xRightarrow{l} \mu''$ such that $(\mu,\mu'') \in \dcal(\scal)$.
\end{definition}

\begin{theorem}[from \cite{Segala1995Compositional}]
    \label{thm:Segala}
    Given two PLTSs $\pcal$ and $\pcal'$, $\pcal \sqsubseteq^t \pcal'$ if and only if there exists a
    probabilistic forward simulation of $\pcal$ by $\pcal'$.
\end{theorem}
\begin{proof}
    The proof is non-trivial and is detailed in Segala's PhD thesis.
    Our approach to this proof is detailed in Appendix \ref{app:ProofSegala}.
\end{proof}

\subsection{Non-probabilistic Systems}

Non-probabilistic systems can be considered as special instances of PLTSs in which all transitions
lead to a Dirac distribution. By identifying each Dirac distribution with the unique state in its
support, we may assume that non-probabilistic systems are modelled by LTSs, which leads to several
simplifications. In particular, probabilistic forward simulation is equivalent to a simpler
simulation relation. On an LTS, we write $q \xRightarrow{l} q'$ if there exists a partial execution
from $q$ to $q'$ whose trace is $l$.

\begin{definition}[Weak Simulation]
    \lean{Weak Simulation}
    %\leanok

    Given two LTSs $\pcal$ and $\pcal'$, a weak simulation of $\pcal$ by $\pcal'$ is a
    relation $\scal \subseteq Q \times Q'$ such that $(q_\star,q_\star') \in \scal$ and,
    for all $(q,q') \in \scal$ and $(q,l,p) \in T$,
    there exists $q' \xRightarrow{l} p'$ such that $(p,p') \in \scal$.
\end{definition}

\begin{lemma}
    \label{lem:weaksim}
    Given two LTSs $\pcal$ and $\pcal'$, $\pcal \sqsubseteq^t \pcal'$ if and only if there exists a
    weak simulation of $\pcal$ by $\pcal'$.
\end{lemma}
\begin{proof}
    According to Theorem \ref{thm:Segala}, it suffices to prove that
    there exists a weak simulation of $\pcal$ by $\pcal'$ if and only if
    there exists a probabilistic forward simulation of $\pcal$ by $\pcal'$. For one direction,
    it suffices to note that any weak simulation between LTSs is a probabilistic forward simulation.
    For the other direction, if $\scal$ is a probabilistic forward simulation of $\pcal$ by $\pcal'$
    then $\{(q,q') \mid \exists (q,\mu) \in \scal, q' \in \bar\mu\}$ is a weak simulation of $\pcal$
    by $\pcal'$.
\end{proof}

There is a natural injection of LTSs into PLTSs, but also a natural projection of PLTSs onto LTSs.
It consists in transforming all probabilistic non-determinism into classical non-determinism.
The projection of a PLTS $\pcal$, denoted $\rho(\pcal)$, is the LTS whose transitions are defined
by $\{(q,l,q') \mid \exists \mu \in \dcal(Q), (q,l,\mu) \in T \wedge q' \in \bar\mu\}$.
$\rho$ is trivially the identity on LTSs, and $\rho$ preserves the set of executions and hence the
set of traces. As a result, the notion of weak simulation can be lifted to PLTSs: $\scal$ is a weak
simulation of $\pcal$ by $\pcal'$ if and only if it is a weak simulation of $\rho(\pcal)$ by
$\rho(\pcal')$. Since $\rho$ is a projection onto LTSs, this is consistent with the previous
definition.

\subsection{Preserving Liveness}

Liveness properties should not be tested with respect to the set of trace distributions and the set 
of traces of a PLTS because these sets are downward-closed and, in particular, they respectively 
always contain the Dirac distribution of the empty trace and the empty trace. Therefore, against 
such sets, liveness properties like termination are not satisfied by any PLTS. This may seem absurd
but it hints at an implicit progress assumption made when evaluating liveness properties. This
progress assumption could be stated as: if a transition is enabled then a transition must fire.
Thus, it discards all executions that are finite and whose last state is not a deadlock. Executions 
satisfying the progress assumptions are said to be maximal and traces of maximal executions are
also called maximal. Non-probabilistic liveness properties should be tested with respect to the set of
maximal traces. Similarly, a scheduler is said to be maximal if it schedules $\bot$ with positive 
probability only from maximal executions and a trace distribution is said maximal if it is induced
by a maximal scheduler. Probabilistic liveness properties should be tested against maximal trace 
distribution.

In the setting of distributed systems suffering from Byzantine failures, another assumption is
necessary in order to define a sensible evaluation of liveness properties. Indeed, since Byzantine
processes behave arbitrarily, they can monopolise the execution preventing correct processes from
taking a step. As a result, a fairness assumption is introduced stating that, in any infinite 
execution, correct processes must be scheduled infinitely often. To model it formally, a fairness 
function $\text{fair} : T \rightarrow \bool$ is introduced which indicates which transitions are 
fair (e.g. transition of correct processes) and which are unfair (e.g. transitions of Byzantine 
processes). However, Byzantine failures also interfere with the notion of deadlock as they are 
allowed to crash at any time, thus, executions may stop if no fair transition is enabled. We call a
fair deadlock any state whose set of enabled transitions is included in the set of unfair 
transition. An execution is said to be fair if (a) it is finite and ends in a fair deadlock, or (b)
it contains infinitely many fair transitions. A trace is said fair if it is induced by a fair 
execution. Non-probabilistic liveness properties will be evaluated with respect to the set of fair 
traces. Similarly, a scheduler is fair if (a) it schedules $\bot$ with positive probability only 
from fair deadlocks and (b) all infinite executions consistent with it are fair. A trace 
distribution is fair if it is induced by a fair scheduler. Probabilistic liveness properties will 
be evaluated against the set of fair trace distributions.

The set of fair traces is denoted $\tcal^f$ and the set of fair trace distributions is denoted 
$\text{tdist}^f(\pcal)$. Unfortunately, inclusion of $\tcal^f$ is not preserved by parallel
composition because synchronisation issues could create new deadlocks. However, in our case, it 
never happens because synchronisation will be used to call sub-processes locally. The modelling 
solution is to use the input-output approach of \cite{InputOutput} where the set of 
labels is partitioned between input labels $In$ and output labels $Out$. We require that the input 
labels are always enabled: $\forall q \in Q, l \in In, \exists \mu \in \dcal(Q), (q,l,\mu) \in T$. 
This requirement is easily fulfilled by adding self-looping transition labelled by each input label
to all states. Moreover, we restrict composition to PLTSs whose outputs labels are distinct. This 
restriction is not problematic either because labels will encode the protocol APIs. Under such 
assumptions, parallel composition does not create new deadlocks and so the inclusion of $\tcal^f$ 
is preserved by composition. As for $\subseteq^t$, inclusion of $\text{tdist}$ is not preserved by 
parallel composition, and the problem cannot be solved by harmless restrictions. Thus, the fair 
trace distribution pre-congruence is defined as: 
\[
    \pcal \sqsubseteq^f \pcal' \iff \forall \pcal'', S \subseteq L, \text{tdist}^f(\pcal \parallel_S \pcal'') \subseteq  \text{tdist}^f(\pcal' \parallel_S \pcal'') 
\]
It enjoys the same algebraic properties as $\sqsubseteq^t$.

Probabilistic forward simulation and weak simulation preserve $\sqsubseteq^t$, but not 
$\sqsubseteq^f$. Looking at the folklore construction made in Lemma \ref{lem:weaksim}, a fair 
execution can be simulated by an unfair execution. If the execution ends with a fair deadlock then
the simulation does not guarantee anything about the last state of the simulating execution.
Moreover, internal fair transitions can be elided, thus, an infinite fair execution can be simulated
by an execution with finitely many fair transitions. We call an execution a fair divergence if it is
fair and its trace is $\tau$. The diagnostic is the following: both simulations do not simulate
fair divergences properly. We will say that a weak simulation $\scal$ is fair if, for all $(q,p) 
\in \scal$, if $q$ fairly diverges then $p$ fairly diverges. However, this simple solution has a
major problem: proving that a state does not fairly diverge is hard. Fortunately, 
\cite{Dongol2022Weak} proposes a solution to a similar problem: equip the simulation with a ranking
function. We adapt their contribution to our setting. 

\begin{definition}[Fair Weak Simulation]
    \lean{FairWeakSimulation}
    %\leanok

    Given two LTSs $\pcal$ and $\pcal'$, a fair weak simulation of $\pcal$ by $\pcal'$ is a 
    relation $\scal \subseteq Q \times Q'$ equipped with a well-founded preorder $\le$ on $Q$ such 
    that $(q_\star,q_\star') \in \scal$ and, for all $(q,q') \in \scal$ and $(q,l,p) \in T$, there 
    exists $q' \xRightarrow{l} p'$ such that $(p,p') \in \scal$. Furthermore, if $l = \tau$ and 
    $q' \xRightarrow{\tau} p'$ contain no fair transition then if $(q,l,p)$ is fair then $q > p$; 
    else $q \ge p$. Moreover, for all $(q,q') \in \scal$, if $q$ is a fair deadlock then $q'$ 
    fairly diverges.
\end{definition}

For probabilistic forward simulation, the problem is more complicated because the construction of
the simulating scheduler is not so clear (see Appendix \ref{app:ProofSegala}). However, the core
idea stays the same even though it is more convoluted.

\begin{definition}[Fair Probabilistic Forward Simulation]
    \lean{FairProbForwardSimulation}
    %\leanok

    Given two LTSs $\pcal$ and $\pcal'$, a fair probabilistic forward simulation of $\pcal$ by 
    $\pcal'$ is a relation $\scal \subseteq Q \times \dcal(Q')$ equipped with a well-founded 
    preorder $\le$ on $Q$ such that $(q_\star,\delta_{q_\star'}) \in \scal$ and, for all $(q,\mu') 
    \in \scal$ and $(q,l,\mu) \in T$, there exists $\mu' \xRightarrow{l} \mu''$ such that 
    $(\mu',\mu'') \in \dcal(\scal)$. Furthermore, let $s$ be the scheduler inducing $\mu' 
    \xRightarrow{\tau} \mu''$ and $w$ the witness of $(\mu',\mu'') \in \dcal(\scal)$, if $l = \tau$
    and there exist $q' \in \bar\mu' \xRightarrow{\tau} p'$ consistent with $s$ which contain no 
    fair transition then, for all $p \in \bar\mu$, if there exists $\mu_p \in \dcal(Q')$ such that
    $w(p,\mu_p) > 0$ and $\mu_p(p') > 0$ then if $(q,l,\mu)$ is fair then $q > p$; else $q \ge p$.
    Moreover, for all $(q,\mu') \in \scal$, if $q$ is a fair deadlock then there exists a fair 
    scheduler $s$ such that $\tfrak_{\mu',s} = \delta_\tau$.
\end{definition}

\begin{theorem}
    \label{thm:Segalafair}
    Given two PLTSs $\pcal$ and $\pcal'$, if there exists a fair probabilistic forward simulation 
    of $\pcal$ by $\pcal'$ then $\pcal \sqsubseteq^f \pcal'$.
\end{theorem}
\begin{proof}
    [TODO]
\end{proof}

\begin{lemma}
    \label{lem:fairweaksim}
    Given two LTSs $\pcal$ and $\pcal'$, if there exists a fair weak simulation of $\pcal$ by 
    $\pcal'$ then $\pcal \sqsubseteq^f \pcal'$.
\end{lemma}
\begin{proof}
    It suffices to note that a fair weak simulation between LTSs is a fair probabilistic forward
    simulation.
\end{proof}

\subsection{Partial Information}

We encode partial information by restrictions on the set of schedulers.
Intuitively, if the scheduler does not have access to all information then it might
not be able to distinguish some states or some labels; thus, some distinct
executions might look the same to it and it should schedule the same
probability distribution of transitions.

\begin{definition}[Adversary]
    An adversary is a pair of observation functions $(\ocal_Q, \ocal_L)$ which
    respectively map states and labels to state observation and label observation
    such that, for all $t,t' \in T$, if $src(t) = src(t')$ and
    $\ocal_L(lbl(t)) = \ocal_L(lbl(t'))$ then $t = t'$.
\end{definition}

The injectivity of $\ocal_L$ on the set of transitions enabled in a given
state ensures that partial information does not prevent the adversary from resolving
non-determinism. An adversary $(\ocal_Q,\ocal_L)$ induces an observation $\ocal_{\ecal}$
on executions defined by
$\ocal_{\ecal}(q_0.l_0\cdots) = \ocal_Q(q_0).\ocal_L(l_0)\cdots$.
A scheduler $s$ is consistent with an adversary if, for all $e,e' \in \ecal^*$,
$\ocal_{\ecal}(e) = \ocal_{\ecal}(e') \implies s(e) = s(e')$. The omniscient adversary
is encoded by setting the observation functions to the identity.

Our simulations are sound for partial-information adversaries because they were
designed to be sound for omniscient adversaries. However, there is a completeness
gap as partial information makes more reductions sound. Our chosen solution is to
build a PLTS on which the omniscient adversary is equivalent to the
partial-information adversary on the original PLTS. This is the belief construction.

\begin{definition}[Belief PLTS]
    Given a PLTS $\pcal$ and an adversary $\acal :=(\ocal_Q,\ocal_L)$, the belief PLTS
    $\pcal_{\acal}$ is defined as the tuple
    $(\{\mu \in \dcal(Q) \mid \ocal_Q(\bar\mu) = \{\cdot\} \},
    \delta_{q_\star}, \ocal_L(L), T_{\acal})$
    where
    \[
        \mu \xrightarrow{l} \mu' \in T_ {\acal} \iff
        \exists \mu \xrightarrow{l} \mu'' \in H,
        \mu' = (q \mapsto
        \frac{\mathbb{1}_{\ocal_Q(q) = o} \cdot \mu''(q)}{\mu''(\{q \mid \ocal_Q(q) = o\})})
        \mapsto \mu''(\{q \mid \ocal_Q(q) = o\})
    \]
\end{definition}

\begin{theorem}
    \label{thm:partialinfosound}
    Given a PLTS $\pcal$ and an adversary $\acal$, the set of trace distributions
    of $\pcal$ induced by a scheduler consistent with $\acal$ is equal to $ \text{tdist}(\pcal)$.
\end{theorem}

When considering the composition of two PLTSs $\pcal$ and $\pcal'$ scheduled by
their respective adversaries $\acal$ and $\acal'$, we need to define a composite
adversary. Given $\pcal$ and $\pcal'$ two PLTSs and $\acal$ and $\acal'$ two adversaries
(one for each PLTS) and $S \subseteq L \cap L'$, we say that $\acal$ and $\acal'$ are
compatible on $S$ if (i), for all $l \in S$, $\ocal_L(l) = \ocal'_L(l)$;
(ii) $\ocal_L(L \setminus S) \cap \ocal_L(S) = \emptyset$ and
(iii) $\ocal'_L(L' \setminus S) \cap \ocal'_L(S) = \emptyset$.
The three requirements state that the adversaries of each PLTS get the same
information from the synchronised labels and distinguish synchronised labels
from the others. If $\acal$ and $\acal'$ are compatible on $S$, we can define
$\acal \parallel_S \acal'$ as
\[
    ((q,q') \mapsto (\ocal_Q(q),\ocal'_Q(q')),
    l \mapsto
    \begin{cases}
        \ocal_L(l) \text{ if } l \in L \\
        \ocal_L'(l) \text{ if } l \in L'
    \end{cases})
\]

\begin{theorem}
    \label{thm:partialinfocompo}
    Given two PLTSs $\pcal$ and $\pcal'$ and their respective adversary $\acal$
    and $\acal'$ and $S \subseteq L\cap L'$, if $\acal$ and $\acal'$ are compatible
    on $S$ then $\pcal_{\acal} \parallel_S \pcal'_{\acal'}$ and
    $(\pcal \parallel_S \pcal')_{\acal \parallel_S \acal'}$ are isomorphic.
\end{theorem}
\begin{proof}
    For clarity, let us denote $\pcal_1 := \pcal_{\acal} \parallel_S \pcal'_{\acal'}$ and
    $\pcal_2 := (\pcal \parallel_S \pcal')_{\acal \parallel_S \acal'}$.
    First, note that $Q_1$ is in $\dcal(Q) \times \dcal(Q')$
    whereas $Q_2$ is in $\dcal(Q \times Q')$ and there is no bijection between
    the two. However, one can prove that $Q_2$ is actually included in
    $\{\mu \in \dcal(Q \times Q') \mid \mu := \mu' \otimes \mu''\}$, thus,
    $h : (\mu,\mu') \mapsto \mu \otimes \mu'$ is a bijection between $Q_1$ and $Q_2$.
    It remains to prove the morphism part which consists in unfolding definitions.
\end{proof}
